\documentclass[CJK]{beamer}
\usepackage{CJKutf8}
\usepackage{beamerthemesplit}
\usetheme{Malmoe}
\useoutertheme[footline=authortitle]{miniframes}
\usepackage{amsmath}
\usepackage{amssymb}
\usepackage{graphicx}
\usepackage{eufrak}
\usepackage{color}
\usepackage{slashed}
\usepackage{simplewick}
\usepackage{tikz}
\usepackage{tcolorbox}
\graphicspath{{../figures/}}
%%figures
\def\lfig#1#2{\includegraphics[width=#1 in]{#2}}
\def\addfig#1#2{\begin{center}\includegraphics[width=#1 in]{#2}\end{center}}
\def\wulian{\includegraphics[width=0.18in]{emoji_wulian.jpg}}
\def\bigwulian{\includegraphics[width=0.35in]{emoji_wulian.jpg}}
\def\bye{\includegraphics[width=0.18in]{emoji_bye.jpg}}
\def\bigbye{\includegraphics[width=0.35in]{emoji_bye.jpg}}
\def\huaixiao{\includegraphics[width=0.18in]{emoji_huaixiao.jpg}}
\def\bighuaixiao{\includegraphics[width=0.35in]{emoji_huaixiao.jpg}}
%% colors
\def\blacktext#1{{\color{black}#1}}
\def\bluetext#1{{\color{blue}#1}}
\def\redtext#1{{\color{red}#1}}
\def\darkbluetext#1{{\color[rgb]{0,0.2,0.6}#1}}
\def\skybluetext#1{{\color[rgb]{0.2,0.7,1.}#1}}
\def\cyantext#1{{\color[rgb]{0.,0.5,0.5}#1}}
\def\greentext#1{{\color[rgb]{0,0.7,0.1}#1}}
\def\darkgray{\color[rgb]{0.2,0.2,0.2}}
\def\lightgray{\color[rgb]{0.6,0.6,0.6}}
\def\gray{\color[rgb]{0.4,0.4,0.4}}
\def\blue{\color{blue}}
\def\red{\color{red}}
\def\green{\color{green}}
\def\darkgreen{\color[rgb]{0,0.4,0.1}}
\def\darkblue{\color[rgb]{0,0.2,0.6}}
\def\skyblue{\color[rgb]{0.2,0.7,1.}}
%%control
\def\be{\begin{equation}}
\def\ee{\nonumber\end{equation}}
\def\bea{\begin{eqnarray}}
\def\eea{\nonumber\end{eqnarray}}
\def\bch{\begin{CJK}{UTF8}{gbsn}}
\def\ech{\end{CJK}}
\def\bitem{\begin{itemize}}
\def\eitem{\end{itemize}}
\def\bcenter{\begin{center}}
\def\ecenter{\end{center}}
\def\bex{\begin{minipage}{0.2\textwidth}\includegraphics[width=0.6in]{jugelizi.png}\end{minipage}\begin{minipage}{0.76\textwidth}}
\def\eex{\end{minipage}}
\def\chtitle#1{\frametitle{\bch#1\ech}}
\def\bmat#1{\left(\begin{array}{#1}}
\def\emat{\end{array}\right)}
\def\bcase#1{\left\{\begin{array}{#1}}
\def\ecase{\end{array}\right.}
\def\bmini#1{\begin{minipage}{#1\textwidth}}
\def\emini{\end{minipage}}
\def\tbox#1{\begin{tcolorbox}#1\end{tcolorbox}}
\def\pfrac#1#2#3{\left(\frac{\partial #1}{\partial #2}\right)_{#3}}
%%symbols
\def\sone{$\star$}
\def\stwo{$\star\star$}
\def\sthree{$\star\star\star$}
\def\sfour{$\star\star\star\star$}
\def\sfive{$\star\star\star\star\star$}
\def\rint{{\int_\leftrightarrow}}
\def\roint{{\oint_\leftrightarrow}}
\def\stdHf{{\textit{\r H}_f}}
\def\deltaH{{\Delta \textit{\r H}}}
\def\ii{{\dot{\imath}}}
\def\skipline{{\vskip0.1in}}
\def\skiplines{{\vskip0.2in}}
\def\lagr{{\mathcal{L}}}
\def\hamil{{\mathcal{H}}}
\def\vecv{{\mathbf{v}}}
\def\vecx{{\mathbf{x}}}
\def\vecy{{\mathbf{y}}}
\def\veck{{\mathbf{k}}}
\def\vecp{{\mathbf{p}}}
\def\vecn{{\mathbf{n}}}
\def\vecA{{\mathbf{A}}}
\def\vecP{{\mathbf{P}}}
\def\vecsigma{{\mathbf{\sigma}}}
\def\hatJn{{\hat{J_\vecn}}}
\def\hatJx{{\hat{J_x}}}
\def\hatJy{{\hat{J_y}}}
\def\hatJz{{\hat{J_z}}}
\def\hatj#1{\hat{J_{#1}}}
\def\hatphi{{\hat{\phi}}}
\def\hatq{{\hat{q}}}
\def\hatpi{{\hat{\pi}}}
\def\vel{\upsilon}
\def\Dint{{\mathcal{D}}}
\def\adag{{\hat{a}^\dagger}}
\def\bdag{{\hat{b}^\dagger}}
\def\cdag{{\hat{c}^\dagger}}
\def\ddag{{\hat{d}^\dagger}}
\def\hata{{\hat{a}}}
\def\hatb{{\hat{b}}}
\def\hatc{{\hat{c}}}
\def\hatd{{\hat{d}}}
\def\hatN{{\hat{N}}}
\def\hatH{{\hat{H}}}
\def\hatp{{\hat{p}}}
\def\Fup{{F^{\mu\nu}}}
\def\Fdown{{F_{\mu\nu}}}
\def\newl{\nonumber \\}
\def\vece{\mathrm{e}}
\def\calM{{\mathcal{M}}}
\def\calT{{\mathcal{T}}}
\def\calR{{\mathcal{R}}}
\def\barpsi{\bar{\psi}}
\def\baru{\bar{u}}
\def\barv{\bar{\upsilon}}
\def\qeq{\stackrel{?}{=}}
\def\torder#1{\mathcal{T}\left(#1\right)}
\def\rorder#1{\mathcal{R}\left(#1\right)}
\def\contr#1#2{\contraction{}{#1}{}{#2}#1#2}
\def\trof#1{\mathrm{Tr}\left(#1\right)}
\def\trace{\mathrm{Tr}}
\def\comm#1{\ \ \ \left(\mathrm{used}\ #1\right)}
\def\tcomm#1{\ \ \ (\text{#1})}
\def\slp{\slashed{p}}
\def\slk{\slashed{k}}
\def\calp{{\mathfrak{p}}}
\def\veccalp{\mathbf{\mathfrak{p}}}
\def\Tthree{T_{\tiny \textcircled{3}}}
\def\pthree{p_{\tiny \textcircled{3}}}
\def\dbar{{\,\mathchar'26\mkern-12mu d}}
\def\erf{\mathrm{erf}}
\def\const{\mathrm{constant}}
\def\pheat{\pfrac p{\ln T}V}
\def\vheat{\pfrac V{\ln T}p}
%%units
\def\fdeg{{^\circ \mathrm{F}}}
\def\cdeg{^\circ \mathrm{C}}
\def\atm{\,\mathrm{atm}}
\def\angstrom{\,\text{\AA}}
\def\SIL{\,\mathrm{L}}
\def\SIkm{\,\mathrm{km}}
\def\SIyr{\,\mathrm{yr}}
\def\SIGyr{\,\mathrm{Gyr}}
\def\SIV{\,\mathrm{V}}
\def\SImV{\,\mathrm{mV}}
\def\SIeV{\,\mathrm{eV}}
\def\SIkeV{\,\mathrm{keV}}
\def\SIMeV{\,\mathrm{MeV}}
\def\SIGeV{\,\mathrm{GeV}}
\def\SIcal{\,\mathrm{cal}}
\def\SIkcal{\,\mathrm{kcal}}
\def\SImol{\,\mathrm{mol}}
\def\SIN{\,\mathrm{N}}
\def\SIHz{\,\mathrm{Hz}}
\def\SIm{\,\mathrm{m}}
\def\SIcm{\,\mathrm{cm}}
\def\SIfm{\,\mathrm{fm}}
\def\SImm{\,\mathrm{mm}}
\def\SInm{\,\mathrm{nm}}
\def\SImum{\,\mathrm{\mu m}}
\def\SIJ{\,\mathrm{J}}
\def\SIW{\,\mathrm{W}}
\def\SIkJ{\,\mathrm{kJ}}
\def\SIs{\,\mathrm{s}}
\def\SIkg{\,\mathrm{kg}}
\def\SIg{\,\mathrm{g}}
\def\SIK{\,\mathrm{K}}
\def\SImmHg{\,\mathrm{mmHg}}
\def\SIPa{\,\mathrm{Pa}}

\def\courseurl{https://github.com/zqhuang/SYSU\_TD}

\def\tpage#1#2{
\begin{frame}
\begin{center}
\begin{Large}
\bch
热学 \\
第#1讲 #2

{\vskip 0.3in}

黄志琦

\ech
\end{Large}
\end{center}

\vskip 0.2in

\bch
教材：《热学》第二版，赵凯华，罗蔚茵，高等教育出版社
\ech

\bch
课件下载
\ech
\courseurl
\end{frame}
}

\title{Lesson 05 - Quantum Statistics}
  \author{}
  \date{}
\begin{document}
\tpage{5}{量子统计}

\section{Reivew}

\begin{frame}
\chtitle{上一讲内容回顾}
\bch
\bitem
\item{能均分定理}
\item{理想摩尔定体热容 = $\frac{R}{2}\times$ 二次型能量的项数}
\item{概率的熵，多可能性事件的熵，多次重复多可能性事件的熵 （理解即可）}
\item{孤立系统的熵增大原理$\Rightarrow$热平衡与热力学温度的数学定义（理解即可）}
\eitem
\ech
\end{frame}


\begin{frame}
\chtitle{本讲内容}
\bch
\bitem
\item{离散能级和实际摩尔定体热容}
\item{玻色气体和费米气体}
\eitem
\ech
\end{frame}

\section{Quantum States}

\begin{frame}
\bch
上节课介绍的能均分定理假设了分子的能量分布是连续的。

\skipline

这对束系统的能量一般不成立。下面我们计算一些简单的束缚系统的量子能级。



\ech
\end{frame}

\begin{frame}
\chtitle{束缚系统量子态的离散性-谐振子}
\bch
{\small 
考虑一个质量为$m$，固有圆频率为$\omega$(即弹性系数为$m\omega^2$)的谐振子
\addfig{1.5}{harmonic_oscillator.png}
现在考虑所有能量小于$E$的谐振子的态的个数。若能量小于$E$，谐振子的位移$x$和动量$\calp$满足
$$\frac{1}{2}m \omega^2 x^2 + \frac{p^2}{2m} \le E $$
写成椭圆方程标准形式：
$$\frac{x^2} {a^2} + \frac{p^2}{b^2} \le 1 $$
其中$a = \frac{1}{\omega}\sqrt{\frac{2E}{m}}$, $b = \sqrt{2mE}$。故上述方程描述的相空间椭圆的面积为$ \pi a b = \frac{2\pi E}{\omega}$。
}
\ech
\end{frame}


\begin{frame}
\chtitle{束缚系统量子态的离散性-谐振子}
\bch
{\small 
谐振子能量小于$E$的相空间的态的个数为
$$n = \frac{2\pi E}{ h \omega}$$
若把态按能量从低到高排，则第$n$个态的能量为
$$ E_n \approx \frac{ n h\omega}{2\pi}$$
量子力学的严格计算结果为
$$ E_n = \frac{ (n+1/2) h\omega}{2\pi}\ \ (n = 0,1,2,\ldots)$$
量子力学里经常用一个符号$\hbar = \frac{h}{2\pi}$，则
$$ E_n = \left(n + \frac{1}{2}\right) \hbar\omega\ \ (n = 0,1,2,\ldots)$$

}
\ech
\end{frame}


\begin{frame}
\chtitle{振动特征温度}
\bch
{\small 
显然，要使处于能量最低($n=0$)态的谐振子发生振动（即到$n>0$的态去)，最起码要吸收$\hbar \omega$的能量。

\skipline

如果$kT\ll \hbar\omega$ ，则谐振子无法被热运动激发（因分子碰撞只能把$\sim kT$的动能传递给谐振子)。我们把$\Theta_{\rm osc} = \frac{\hbar\omega}{k}$叫做振动特征温度。
}
\ech
\end{frame}


\begin{frame}
\chtitle{振动特征温度的数量级}
\bch
{\small
我们仍然考虑一个双原子气体分子，假设是共价键，典型的键能为$100\SIkcal/\SImol$，即每个键结合能为
$E_b\sim \frac{4.2\times 10^5J}{6.02\times 10^{23}} \sim 0.7\times 10^{-18}J$

我们近似取弹性系数
$m\omega^2 \sim  \frac{2E_b}{(1 \angstrom)^2} \sim 140 \SIN/\SIm$，并取等效质量$m = \frac{A_1A_2}{(A_1+A_2)\times 6.02\times 10^{26}} \SIkg$，则$\omega \sim 10^{14}\SIHz \times \sqrt{\frac{8.4(A_1+A_2)}{A_1A_2}} $

故振动特征温度
$$\Theta_{\rm osc} =\frac{\hbar \omega}{k} \sim 10^3\SIK \times \sqrt{\frac{5(A_1+A_2)}{A_1A_2}} $$
该估算比较粗糙。实际的双原子气体的$\Theta_{\rm osc}$一般在若干个$10^3\SIK$之上。所以{\blue 室温下大多数双原子气体的振动自由度不被激发。}

\skipline

{\blue 固体如果有比较强的键，则有些振动自由度也会被冻结，出现$C^{\rm mol}_V \ll 3R$的情况。}
}
\ech
\end{frame}


\begin{frame}
\chtitle{转动的离散能级}
\bch
现在考虑绕固定轴的逆时针转动，设转动惯量为$I$。广义位置为转动角$\theta$，广义动量为$\calp = I\omega$ (这里$\omega = \frac{d\theta}{dt}$为转动角速度）。

若要转动能量低于$E$，则$\frac{1}{2}I\omega^2 \le E$，或者写成$0\le \calp \le \sqrt{2I E}$。又$0\le \theta \le 2\pi$。则相空间内能量低于$E$的态个数为
$$n = \frac{2\pi \sqrt{2IE}}{h} = \frac{\sqrt{2IE}}{\hbar}$$
若把态按能量从低到高排，则第$n$个态的能量为
$$E_n \approx \frac{n^2\hbar^2}{2I}$$
严格的量子力学计算结果为
$$E_n = \frac{n(n+1)\hbar^2}{2I}$$
\ech
\end{frame}


\begin{frame}
\chtitle{转动特征温度}
\bch
同样道理，若$kT\ll \frac{\hbar^2}{2I}$，则热运动无法激发转动，我们把$\Theta_{\rm rot} = \frac{\hbar^2}{2kI}$叫做转动特征温度。

\ech
\end{frame}


\begin{frame}
\chtitle{转动特征温度的数量级}
\bch
我们考虑一个双原子气体分子，设两个原子间的距离为$r$，两个原子的原子量分别为$A_1$和$A_2$，则转动惯量为
$$I = 10^{-46} \SIkg\cdot\SIm^2 \times \frac{ A_1A_2}{6.02(A_1+A_2)} \left(\frac{r}{1\angstrom}\right)^2 $$
 特征温度为
$$\Theta_{\rm rot} =  \frac{\hbar^2}{2kI} = 24\SIK \times \frac{(A_1+A_2)}{ A_1A_2} \left(\frac{1\angstrom}{r}\right)^2$$
一般都远远低于室温（特别是当$A_1$, $A_2$都较大时）。

所以{\blue 室温下气体分子的转动自由度一般是被激发的}。
\ech
\end{frame}


\begin{frame}
\chtitle{思考题}
\bch
先考虑简单的一维情形，设分子被限制在$0<x<L$的一维盒子里。
\bitem
\item{能量小于$E$的相空间的态有多少个？}
\item{估算第$n$个能级的能量。}
\eitem
\ech
\end{frame}

\begin{frame}
\chtitle{平动特征温度}
\bch
把思考题的结果推广到三维，就可以得到边长为$L$的容器内气体的平动特征温度：
$$\Theta_{\rm trans} = \frac{h^2}{8mkL^2} = 2.4\times 10^{-18}\SIK \times \frac{1}{A}\left(\frac{1\SIm}{L}\right)^2$$
可见，在{\blue 宏观尺度的容器里，平动自由度一般都是被激发的}。
\ech
\end{frame}


\begin{frame}
\chtitle{双原子气体的摩尔定体热容随温度的变化}
\bch
如果从低温开始到高温，双原子气体的转动自由度，振动自由度会被相继激发。摩尔定体热容就从$\frac{3}{2}R$增长到$\frac{5}{2}R$再到$\frac{7}{2}R$.

\addfig{2.5}{CVmol_T.png}

请再参考教材85页 图2-14
\ech
\end{frame}

\section{Bosons and Fermions}

\begin{frame}
\bch
需要理解的内容讲完了。

\skiplines

下面我要开始乱讲，听不懂不要慌，因为——不——考。
\ech
\end{frame}

\begin{frame}
\chtitle{玻尔兹曼不知道的绕口令}
\bch
在玻尔兹曼的时代，量子力学还没被发现，遑论后来量子力学里的再次革命——“二次量子化”。

\skiplines

在二次量子化之后，我们把“测量到a态上有两个粒子”这句话重新进行了描述：“测量到有两个a粒子的态”。

\skipline

\wulian绕口令说得不错
\ech
\end{frame}


\begin{frame}
\chtitle{玻尔兹曼的描述}
\bch
\addfig{1.2}{E_classic.png}

测到a态有3个光子，b态有2个光子，c态有1个光子。
\ech
\end{frame}


\begin{frame}
\chtitle{二次量子化描述}
\bch
\bmini{0.41}
\addfig{1.5}{E_quantum.png}
\emini
\bmini{0.55}
测到光子场处在(3a，2b, 1c)态。

\skiplines

我们以前描述的“粒子态”是用来描述单个粒子的（可以说某粒子处在某态）。而二次量子化之后的“场态”则是描述比粒子更基本的“场”的（\bye 你们已经穿越到研究生噩梦课程《量子场论》了）。

\emini

\ech
\end{frame}


\begin{frame}
\chtitle{玻色子和费米子}
\bch
根据二次量子化的对易假设不同，可以得到两种性质不同的粒子。

\bitem
\item{可以在一个“粒子态”上放任意个粒子，我们称之为玻色子(Boson)。例如光子。}
\item{最多只能在一个“粒子态”上放一个粒子，我们称之为费米子(Fermion)。例如电子，质子。}
\eitem

两个费米子不能在完全相同的态上，这称为“泡利(Pauli)不相容原理”。

注意粒子往往有一些内禀自由度（自旋等）。“泡利不相容原理”仅对“态”
\ech
\end{frame}


\begin{frame}
\chtitle{费米场的态}
\bch
\addfig{1.5}{E_fermion.png}
\ech
\end{frame}


\begin{frame}
\chtitle{稍作修改的“万能法则”}
\bch
我们现在要对“万能法则”稍作修改：场态$(n_a a, n_bb, n_c c,\ldots)$出现的概率正比于
$$e^{-\frac{n_a(\varepsilon_a-\mu)+n_b(\varepsilon_b-\mu)+n_c(\varepsilon_c-\mu)+\ldots}{kT}} = e^{-\frac{\varepsilon_{\rm total} - \mu_{\rm total}}{kT}}$$

\skiplines

这里的$\mu$开始变得重要了，它称为化学势，是一个依赖于温度和总粒子数的量。它随着温度升高而降低，随着总粒子数增大而增大。
\ech
\end{frame}


\begin{frame}
\chtitle{费米-狄拉克分布(Fermi-Dirac distribution)}
\bch
现在考虑某粒子态a，对费米子而言，粒子态a上可以有$n=0,1$个粒子，出现概率依次正比于$e^{-\frac{n(\varepsilon_a-\mu)}{kT}}$。
平均粒子数为

\bea
\bar{n} &=& \frac{\sum_{n=0}^1 ne^{-\frac{n(\varepsilon_a-\mu)}{kT}}}{\sum_{n=0}^1 e^{-\frac{n(\varepsilon_a-\mu)}{kT}}}  \newl
&=& \frac{ e^{-\frac{\varepsilon_a-\mu}{kT}}}{1+ e^{-\frac{\varepsilon_a-\mu}{kT}}}  \newl
&=& \frac{ 1}{e^{\frac{\varepsilon_a-\mu}{kT}}+1}  
\eea

\ech
\end{frame}


\begin{frame}
\chtitle{玻色-爱因斯坦分布(Bose-Einstein distribution)}
\bch
玻色子的情况要稍复杂，对$n$的求和要从$0$到$\infty$。
{\scriptsize
\bea
\bar{n} &=& \frac{\sum_{n=0}^{\infty} ne^{-\frac{n(\varepsilon_a-\mu)}{kT}}}{\sum_{n=0}^1 e^{-\frac{n(\varepsilon_a-\mu)}{kT}}}  \newl
&=& \left. \frac{\sum_{n=0}^{\infty} ne^{-nx}}{\sum_{n=0}^1 e^{-nx}}\right\vert_{x = \frac{\varepsilon_a-\mu}{kT}}  \newl
&=& \left. \frac{-\frac{d}{d x}\sum_{n=0}^{\infty} e^{-nx}}{\sum_{n=0}^1 e^{-nx}}\right\vert_{x = \frac{\varepsilon_a-\mu}{kT}}  \newl
&=& \left. \frac{-\frac{d}{d x}\frac{1}{1- e^{-x}} }{ \frac{1}{1- e^{-x}} }\right\vert_{x = \frac{\varepsilon_a-\mu}{kT}}  \newl
&=& \frac{ 1}{e^{\frac{\varepsilon_a-\mu}{kT}}-1}  
\eea
显然，对玻色子，$\mu$不能高于基态（能量最低的态，注意这可以是多个能量相等的态）的能量。
}

\ech
\end{frame}

\begin{frame}
\chtitle{$T\rightarrow 0K$时的极限}
\bch
当温度$T\rightarrow 0$，不管是费米子还是玻色子，都尽量往低能级跑。于是产生了两种结果：
\bitem
\item{玻色子都集中在基态，形成一个“凝聚体”，称为“{\blue 玻色-爱因斯坦凝聚}”（BE condensation)。}
\item{费米子则因为每个态最多只能放一个粒子，就只能按能量从低到高选择不同的态依次填充。当所有费米子都用完时，填到的那个态的能量$\varepsilon_F$称为{\blue 费米能}。}
\eitem

\skipline

费米气体在$0\SIK$时也会有大量粒子处在激烈运动的非基态，形成的动理压强称为{\blue 简并压}(degeneracy pressure)。这完全是一个量子力学效应。
\ech
\end{frame}


\begin{frame}
\chtitle{思考题}
\bch
试证明当$T\rightarrow 0$时
\bitem
\item{玻色子$\mu \rightarrow $基态能量}
\item{而费米子$\mu \rightarrow \varepsilon_F$}
\eitem

\ech
\end{frame}

\begin{frame}
\chtitle{思考题：简并压和费米能的关系}
\bch
考虑仅有平动自由度的费米气体: $\varepsilon = \frac{\calp^2}{2m}$。设粒子数密度为$n$。

\skiplines

回忆气体动理压强公式：$p =\frac{2}{3}n\bar\varepsilon$ (非相对论情形）及$p =\frac{1}{3}n\bar\varepsilon$ （相对论情形）。对费米气体，证明简并压
\bitem
\item{在非相对论的情形为$\frac{2}{5}n\varepsilon_F$}
\item{在相对论的情形为$\frac{1}{4}n\varepsilon_F$}
\eitem
\ech
\end{frame}


\begin{frame}
\chtitle{估算费米能}
\bch
{\small
考虑某个内禀态数为$g$(例如电子有自旋向上和自旋向下两个内禀态）的费米气体，当温度降为$0\SIK$，动量低于$\calp_{\max}$的态的总数为
$$ N =  g \frac{V \left(\frac{4\pi}{3} \calp_{\max}^3\right)}{h^3} $$
由于$0\SIK$时每个态都填了一个费米子，$n = N/V$即是粒子数密度，我们得到

非相对论情形
$$\varepsilon_F = \frac{\calp_{\max}^2}{2m} =  \left(\frac{3n}{4\pi g}\right)^{2/3} \frac{h^2}{2m} = \left(\frac{6\pi^2 n}{ g}\right)^{2/3} \frac{\hbar^2}{2m}$$

极端相对论情形
$$\varepsilon_F = c\calp_{\max} =  \left(\frac{3n}{4\pi g}\right)^{1/3} hc = \left(\frac{6\pi^2 n}{ g}\right)^{1/3} \hbar c$$
}
\ech
\end{frame}


\begin{frame}
\chtitle{解除量子特性需要的温度}
\bch
现在考虑从$0\SIK$逐渐升高温度。
\bitem
\item{玻色子：逐渐有粒子逃离基态。到某个临界温度时$\Theta_B$时，大多数玻色子逃离基态，BE凝聚被破坏。这个临界温度称为{\blue 玻色-爱因斯坦凝聚点}。}
\item{费米子：当从$0\SIK$逐渐升高温度，开始有费米子跑到能量高于$\varepsilon_F$的态去。当$T> \Theta_F = \varepsilon_F/k$时，可以认为费米气体的简并压解除。这个温度$\Theta_F$称为{\blue 费米简并温度}。}
\eitem
\ech
\end{frame}

\begin{frame}
\chtitle{金属内电子的简并温度和简并压}
\bch
金属中电子的数密度和原子数密度为同一数量级$n\sim 10^{29} \SIm^{-3}$ （还记得估算金属原子平均间距为$2\angstrom$左右吗)。
电子质量$m\approx 10^{-30} \SIkg$。电子内禀自由度为$g = 2$。由此估算出
$$\varepsilon_F = \left(\frac{6\pi^2 n}{ g}\right)^{2/3} \frac{\hbar^2}{2m} \sim 10^{-18}J $$
从而简并温度
$$\Theta_F \sim 8\times 10^4 \SIK$$


所以金属的电子在室温下有简并压。我们还可以估算出简并压为
$$ p = \frac{2}{5} n \varepsilon_F \sim 4\times 10^{10}\SIPa$$


金属抗压几乎全靠的就是电子简并压。参见教材101页表2-10。

\ech
\end{frame}


\begin{frame}
\chtitle{白矮星和中子星}
\bch
\bitem
\item{白矮星是靠电子简并压抗衡引力的高密度天体。}
\item{中子星是靠中子简并压抗衡引力的高密度天体。}
\eitem

课本上有些定性的讨论，参考102-106页。我们不再展开。

\ech
\end{frame}

\begin{frame}
\chtitle{光子气体}
\bch
光子的能量$\varepsilon = \calp c$，基态能量为$0$。当有一个温度为$T$的辐射源和光子气体达到热平衡时，光子的化学势达到$\mu = 0$。
可以计算光子气体的能量为
$$\int_0^\infty \frac{\calp c}{e^{\frac{\calp c}{kT}}-1} \frac{g V 4\pi\calp^2 d\calp}{h^3} $$
其中$g=2$是光子的内禀态数（自旋为$1$和$-1$)，$V$是体积。

\ech
\end{frame}


\begin{frame}
\chtitle{第五周作业(序号接第四周）}
\bch
\bitem
\item[13]{}
\item[14]{}
\item[15]{}
\eitem
\ech
\end{frame}

\end{document}
